\begin{abstract}
In 1960 Widrow and Hoff stored the weights of their adaptive neuron in memistors, three-terminal electrochemical cells; their
modern descendants are lithium-ion, metal-oxide and protonic ECRAM. We ask, in simulation, how close such devices come to digital
training when every weight update is made of device pulses, while the forward and backward reads are kept ideal (fp32). We compare
plain SGD and three Tiki-Taka-family algorithms on five device specifications (two AIHWKIT fits mapped one to one, one measured
array, a protonic-inspired scenario and an optimistic historical scenario), on MNIST, a CIFAR-10 subset and a byte-level language
model, with five seeds that share initialisation and data order across arms. Against SGD with exact updates at the same learning
rate, \RVAllThreeFpN{} pairings (Li-ECRAM and protonic-inspired, each with chopped TTv2 or AGAD) stay within the pre-set margins
on all three tasks, and plain SGD exceeds them on every published device. Against a separately tuned digital Adam no pairing
does: on the language model the best analog result is \RVLmAdamBestGap\,\% worse in bits per byte. After per-arm tuning we detect
no growth of the gap with width over a factor of four, without establishing equivalence. A sensitivity sheet against absolute
accuracy lines, with three seeds, establishes only a few one-at-a-time ranges (on CIFAR with chopped TTv2, step asymmetry up to
\RVSpecCttvTwoCifarAsymThr{} and reference error up to \RVSpecCttvTwoCifarRefErrThr), and shows that an earlier asymmetry limit
was mostly lost weight range. A conditional
energy projection for the weight-matrix operations, borrowing the read efficiency of an analog inference chip, gives about
\RVenSysMin--\RVenChipMax$\times$ less energy per weight-row than an H100, reported only where the quality is reached. This
version repairs the simulator after an independent review (matched random streams, a physically applied reset, the device presets
mapped one to one) and reports how much each repair moved the original results. We propose, but did not build, a bench
replication protocol of the 1960 memistor.
\end{abstract}
